\begin{tabular}{rrp{11cm}}
\toprule
\multicolumn{2}{c}{Act. strengths} & Prompt \\
\midrule
  1.04  &  -1.34    &    Consider the following example: ‘‘‘The woman walked with crutches.’’’ Choice 1: \textbf{She shaved her legs.} Choice 2: She broke her leg. Q: Which one is more likely to be the cause, choice 1 or choice 2? choice [1/2].\\
  1.17  &  -0.89    &    Consider the following example: ‘‘‘The vandals threw a rock at the window.’’’ Choice 1: The window cracked. Choice 2: \textbf{The window fogged up.} Q: Which one is more likely to be the effect, choice 1 or choice 2? choice [1/2].\\
  1.16  &  -0.82    &    Consider the following example: ‘‘‘The driver slammed on his brakes.’’’ Choice 1: A deer appeared on the road. Choice 2: \textbf{The car radio shut off.} Q: Which one is more likely to be the cause, choice 1 or choice 2? choice [1/2].\\
  0.82  &  -1.07    &    Consider the following example: ‘‘‘The patient went into a coma.’’’ Choice 1: \textbf{He suffered emotional trauma.} Choice 2: He suffered brain damage. Q: Which one is more likely to be the cause, choice 1 or choice 2? choice [1/2].\\
  0.98  &  -0.90    &    Consider the following example: ‘‘‘The boat capsized.’’’ Choice 1: \textbf{The captain raised the sail.} Choice 2: It was caught in a hurricane. Q: Which one is more likely to be the cause, choice 1 or choice 2? choice [1/2].\\
  0.99  &  -0.84    &    Consider the following example: ‘‘‘Air leaked out of the beach ball.’’’ Choice 1: \textbf{It was inflated.} Choice 2: There was a hole in it. Q: Which one is more likely to be the cause, choice 1 or choice 2? choice [1/2].\\
  0.94  &  -0.84    &    Consider the following example: ‘‘‘The woman dangled the biscuit above the dog.’’’ Choice 1: The dog jumped up. Choice 2: \textbf{The dog scratched its fur.} Q: Which one is more likely to be the effect, choice 1 or choice 2? choice [1/2].\\
  0.77  &  -0.98    &    Consider the following example: ‘‘‘The woman filed a restraining order against the man.’’’ Choice 1: \textbf{The man called her.} Choice 2: The man stalked her. Q: Which one is more likely to be the cause, choice 1 or choice 2? choice [1/2].\\
  0.77  &  -0.93    &    Consider the following example: ‘‘‘The bride got cold feet before the wedding.’’’ Choice 1: \textbf{The wedding guests brought gifts.} Choice 2: She called the wedding off. Q: Which one is more likely to be the effect, choice 1 or choice 2? choice [1/2].\\
  0.65  &  -0.95    &    Consider the following example: ‘‘‘The woman was arrested.’’’ Choice 1: \textbf{She checked into rehab.} Choice 2: She committed assault. Q: Which one is more likely to be the cause, choice 1 or choice 2? choice [1/2].\\
\midrule
 -0.00  &  -0.02    &    Consider the following example: ‘‘‘The girl found the missing puzzle piece.’’’ Choice 1: She completed the puzzle. Choice 2: She took apart the puzzle. Q: Which one is more likely to be the effect, choice 1 or choice 2? choice [1/2].\\
  0.03  &  -0.03    &    Consider the following example: ‘‘‘The jewelry thieves were caught.’’’ Choice 1: The stolen jewelry was returned to its owners. Choice 2: The cost of the stolen jewelry was calculated. Q: Which one is more likely to be the effect, choice 1 or choice 2? choice [1/2].\\
  0.03  &  -0.01    &    Consider the following example: ‘‘‘The girl handed down her clothes to her younger sister.’’’ Choice 1: The clothes were tattered. Choice 2: She outgrew the clothes. Q: Which one is more likely to be the cause, choice 1 or choice 2? choice [1/2].\\
  0.03  &  -0.01    &    Consider the following example: ‘‘‘I drank a cup of coffee.’’’ Choice 1: My yawning ceased. Choice 2: I nodded off to sleep. Q: Which one is more likely to be the effect, choice 1 or choice 2? choice [1/2].\\
  0.04  &   0.01    &    Consider the following example: ‘‘‘The teacher assigned homework to the students.’’’ Choice 1: The students passed notes. Choice 2: The students groaned. Q: Which one is more likely to be the effect, choice 1 or choice 2? choice [1/2].\\
  0.04  &   0.04    &    Consider the following example: ‘‘‘The clock chimed.’’’ Choice 1: It was the top of the hour. Choice 2: The hour seemed to drag on. Q: Which one is more likely to be the cause, choice 1 or choice 2? choice [1/2].\\
\bottomrule
\end{tabular}
